\section{Regulatory phenotype and context}

\subsection{Network calibration and developmental experience}

Network calibration denotes measurable development of synaptic organization, release dynamics, sensory responses, and regulation under perturbation. It must be decomposed into outcomes such as evoked responses, recovery time, variability, and learning rather than inferred from a global impression of “high gain.”

Experimental rat work connects maternal care with epigenetic regulation of glucocorticoid-receptor-related stress processes \cite{ref14}. Human postmortem work associates childhood abuse history with hippocampal glucocorticoid receptor regulation \cite{ref15}. Species, tissue, developmental window, ascertainment, and causal design differ. These studies support developmental embedding as a domain of inquiry; they do not establish that the same epigenetic modification mediates the present PRI–CaV proposal or that acquired states are transmitted through the human germline.

Human stress-methylation evidence is mixed rather than a settled inheritance mechanism. In 113 mother–infant dyads, maternal maltreatment associations did not directly appear in newborn methylation at the assayed sites \cite{ref25}. A separate FKBP5 association replicated in Holocaust offspring, but the transmission route remained unresolved \cite{ref26}. In a population sample of 3,965 adults, previously reported childhood-maltreatment interactions at tested FKBP5 sites were not replicated \cite{ref34}. These studies differ in tissue, timing, cell mixture, genotype and ascertainment. They motivate prospective measurement, not an assertion that trauma is inherited as a specific methylation mark.

Autonomic variation, sleep, pain, sensory load, learned threat, and social reinforcement can change current attention without requiring an early molecular lesion. Adult autistic camouflaging has a validated self-report measure \cite{ref32}, but the broader cascade’s automatic or homeostatic masking remains a proposed transfer across phenotypes. These are competing and interacting pathways. A useful longitudinal analysis must allow compensation, recovery, nonlinear responses, and different trajectories leading to similar observed phenotypes.


\subsection{Regulatory Closure Bias: construct and operationalization}

Regulatory Coupling Orientation names the hypothesis that the relative contribution of external interaction and internally sustained stabilization varies by person, task, and state. Both are present in every person. RCB is the proposed quantitative index:

\begin{equation}
\mathrm{RCB}_{i,\tau,s}=\ln\!\left(\frac{G_{\mathrm{exo}}}{G_{\mathrm{endo}}}\right),
\qquad G_{\mathrm{exo}},G_{\mathrm{endo}}>0 .
\end{equation}

The logarithm is defined only for positive, commensurate gains. These terms must not be assigned from a diagnosis or questionnaire label. An event-rate meta-analysis found disproportionate slowing of ADHD groups under slow Go/No-Go event rates, larger commission-error effects under fast rates, and persisting general group differences; reaction-time variability did not show the predicted event-rate effect \cite{ref27}. That is a bounded task result, not a validated RCB index. One candidate protocol randomizes participants across externally responsive and internally paced conditions with matched difficulty, reward, sensory load, practice, and time on task. Predefine a regulatory error L, such as a weighted, standardized combination of performance variability and recovery following a task perturbation. Estimate condition-specific reductions relative to a common reference.

If an estimated gain is nonpositive or indistinguishable from zero, report the component estimates and their uncertainty; do not force a ratio by adding an arbitrary constant. A prespecified stabilized index may be evaluated as a separate measure but must disclose its dependence on the chosen constant. Repeated sessions should test reliability, construct validity, and incremental prediction beyond arousal, executive function, reward sensitivity, and symptom severity.

The website's positive-gain sliders illustrate the mathematics only. Their values are arbitrary teaching inputs, not normative scores, biological measurements, diagnoses, or estimated coefficients. The main empirical prediction concerns the interaction of independently measured RCB with environment; a favorable response cannot both define and validate the construct in the same data.


\subsection{Hunter/Farmer and Dąbrowski interpretations}

Hunter/Farmer terminology describes ecological renderings of behavior, not ancestral castes or fixed human types. An online foraging study linked attention-deficit traits with exploratory choices \cite{ref16}. In a different Hadza actigraphy study, one or more group members were scored awake or in light sleep in 99.8\% of sampled nighttime epochs across 20 days \cite{ref28}. This illustrates group-level chronotype heterogeneity, but did not measure ADHD, ancestral subtypes, or modern work performance. Such findings motivate task-dependent hypotheses; they do not prove an evolutionary origin of ADHD or universal advantage under uncertainty.

Dąbrowski's framework contributes a vocabulary for instability followed by reorganization \cite{ref17}. Psychomotor, sensory, emotional, intellectual, and imaginational overexcitabilities are retained as historical conceptual dimensions. They are not demonstrated consequences of RCCX or \ensuremath{\alpha_2\delta} biology. “Bifurcation” is used as a model of divergent trajectories, not an empirically identified mathematical phase transition.

The studyable question is whether particular combinations of sensitivity, context, support, and self-regulation predict improved functioning or sustained impairment. Reorganization can be multidirectional, reversible, and domain-specific. Suffering is neither necessary nor sufficient for insight; disability, originality, and achievement must be measured separately. Conversion and dissemination depend on resources and social opportunity as well as individual capacity.


\subsection{Pharmacological retrodiction: provenance, not diagnosis}

Adult reports of broad state changes following a compound can motivate a candidate mechanism. Zvejniece et al. provide experimental evidence for R-phenibut binding to \ensuremath{\alpha_2\delta} alongside its GABA-B pharmacology \cite{ref18}. This offers a reason to investigate both targets, not a way to infer which target produced a particular person's experience.

Retrodiction from an adult response to a developmental cause is underdetermined by multiple targets, kinetics, current state, expectation, and downstream network effects. Similarity of benefit does not identify a shared lesion; differences between related compounds do not uniquely isolate a target. No adult self-report is counted as evidence for PRI or \ensuremath{\alpha_2\delta} developmental pathology here. Pharmacological provenance is recorded separately from the causal ledger, without a treatment proposal or dosing protocol.
